\clearpage
\section{Comparison of \RunPeriod{} results to \PrevRunPeriod{}}
\label{sec:comparison}

This chapter compares the results for the current \RunPeriod data set to the results obtained for the previous \PrevRunPeriod data set. Comparisons of \RunPeriod data to Asimov A sensitivities can be found in \cref{app:data_asimov_comparison}, whilst comparisons between sensitivity studies from the respective periods can be found in \cref{app:asimov_comparison}. Comparisons of data fits and sensitivity studies between the different fitting groups can be found in \cref{app:group_comparison}.

\subsection{Results for \appearance}

% TODO: Update this for run 10
In \cref{fig:ssqth13_vs_dcp_data_compare} the result for \appearance with reactor constraint and T2K only is shown. %The principle difference in the contour with respect to \PrevRunPeriod is seen with the reactor constraint applied, where the new constraint has shifted the best-fit towards a lower value of \ssqthonethree and substantially shrunk the contour along this dimension. The fit without the reactor constraint applied shows a small reduction in the size of the contour since \PrevRunPeriod.

\begin{figure*}[htb]
	\center
	\begin{subfigure}[b]{0.49\textwidth}
		\includegraphics[width=0.9\textwidth]{images/data/comparison/appearance_wrc_normal_1-9d_v_1-9c}
		\caption{\appearance NH with reactor constraint}
	\end{subfigure}
	\begin{subfigure}[b]{0.49\textwidth}
		\includegraphics[width=0.9\textwidth]{images/data/comparison/appearance_wrc_inverted_1-9d_v_1-9c}
		\caption{\appearance IH with reactor constraint}
	\end{subfigure}
	\begin{subfigure}[b]{0.49\textwidth}
		\includegraphics[width=0.9\textwidth]{images/data/comparison/appearance_worc_normal_1-9d_v_1-9c}
		\caption{\appearance NH without reactor constraint}
	\end{subfigure}
	\begin{subfigure}[b]{0.49\textwidth}
		\includegraphics[width=0.9\textwidth]{images/data/comparison/appearance_worc_inverted_1-9d_v_1-9c}
		\caption{\appearance IH without reactor constraint}
	\end{subfigure}
	\caption {
		Contours at 68\% and 90\% CL for \appearance with and without reactor constraint. Normal and inverted hierarchy contours are independent. All 5 samples were used to produce these fits.
	}
	\label{fig:ssqth13_vs_dcp_data_compare}
\end{figure*}
\clearpage

\subsection{Results for \disappearance}

% TODO: Update this for run 10
The result for \disappearance with reactor constraint and T2K only is shown in fig.~\ref{fig:ssqth23_vs_dmsq_data_compare}. %The notable changes in the disappearance fits are a shift of the contour towards higher values of \dmsq (for both normal and inverted hierarchy), and a reduction in the size of the contour, particularly in the \ssqthtwothree dimension.

\begin{figure*}[htb]
	\center
	\begin{subfigure}[b]{0.49\textwidth}
		\includegraphics[width=0.9\textwidth]{images/data/comparison/disappearance_wrc_normal_1-9d_v_1-9c}
		\caption{\disappearance NH with reactor constraint}
	\end{subfigure}
	\begin{subfigure}[b]{0.49\textwidth}
		\includegraphics[width=0.9\textwidth]{images/data/comparison/disappearance_wrc_inverted_1-9d_v_1-9c}
		\caption{\disappearance IH with reactor constraint}
	\end{subfigure}
	\begin{subfigure}[b]{0.49\textwidth}
		\includegraphics[width=0.9\textwidth]{images/data/comparison/disappearance_worc_normal_1-9d_v_1-9c}
		\caption{\disappearance NH without reactor constraint}
	\end{subfigure}
	\begin{subfigure}[b]{0.49\textwidth}
		\includegraphics[width=0.9\textwidth]{images/data/comparison/disappearance_worc_inverted_1-9d_v_1-9c}
		\caption{\disappearance IH without reactor constraint}
	\end{subfigure}
	\caption {
		Contours at 68\% and 90\% CL for \disappearance with and without reactor constraint. Normal and inverted hierarchy contours are independent. All 5 samples were used to produce these fits.
	}
	\label{fig:ssqth23_vs_dmsq_data_compare}
\end{figure*}

%\iffalse
% TODO: Update this for run 10
%A potential explanation for the broadening of the contour, despite the increase in statistics could come from a shift in the value of \ssqthtwothree away from maximal disappearance ($\sin^2\theta_{23} \simeq 0.51$).
%\fi

\clearpage

\subsection{Results for \deltacp}

% TODO: Update this for run 10
The result for \deltacp with reactor constraint and T2K only is shown in fig.~\ref{fig:dcp_data_compare}. %With the reactor constraint applied we observe a small shift away from maximal CP-violation relative to \PrevRunPeriod, with slightly stronger constraints evident.

\begin{figure*}[htb]
	\center
	\begin{subfigure}[b]{0.49\textwidth}
		\includegraphics[width=0.9\textwidth]{images/data/comparison/deltacp_wrc_normal_1-9d_v_1-9c}
		\caption{\deltacp NH with reactor constraint}
	\end{subfigure}
	\begin{subfigure}[b]{0.49\textwidth}
		\includegraphics[width=0.9\textwidth]{images/data/comparison/deltacp_wrc_inverted_1-9d_v_1-9c}
		\caption{\deltacp IH with reactor constraint}
	\end{subfigure}
	\begin{subfigure}[b]{0.49\textwidth}
		\includegraphics[width=0.9\textwidth]{images/data/comparison/deltacp_worc_normal_1-9d_v_1-9c}
		\caption{\deltacp NH without reactor constraint}
	\end{subfigure}
	\begin{subfigure}[b]{0.49\textwidth}
		\includegraphics[width=0.9\textwidth]{images/data/comparison/deltacp_worc_inverted_1-9d_v_1-9c}
		\caption{\deltacp IH without reactor constraint}
	\end{subfigure}
	\caption {
		The expected $\Delta\chi^2$ distribution as a function of \deltacp with and without reactor constraint. Normal and inverted hierarchy contours are independent. All 5 samples were used to produce these fits.
	}
	\label{fig:dcp_data_compare}
\end{figure*}
\clearpage

\subsection{Results for \ssqthonethree}

% TODO: Update this for run 10
The result for \ssqthonethree (without reactor constraint) is shown in fig.~\ref{fig:th13_data_compare}. %A small shift in the best-fit value can be observed, along with a tightening of the contour.

\begin{figure*}[htb]
	\center
	\begin{subfigure}[b]{0.49\textwidth}
		\includegraphics[width=0.9\textwidth]{images/data/comparison/theta13_worc_normal_1-9d_v_1-9c}
		\caption{\ssqthonethree NH without reactor constraint}
	\end{subfigure}
	\begin{subfigure}[b]{0.49\textwidth}
		\includegraphics[width=0.9\textwidth]{images/data/comparison/theta13_worc_inverted_1-9d_v_1-9c}
		\caption{\ssqthonethree IH without reactor constraint}
	\end{subfigure}
	\caption {
		The expected $\Delta\chi^2$ distribution as a function of \ssqthonethree without reactor constraint. Normal and inverted hierarchy contours are independent. All 5 samples were used to produce these fits.
	}
	\label{fig:th13_data_compare}
\end{figure*}

\clearpage

\subsection{Results for \ssqthtwothree}

% TODO: Update this for run 10
The result for \ssqthtwothree with reactor constraint and T2K only is shown in fig.~\ref{fig:th23_data_compare}. %As noted in the 2D fit, the contour is somewhat narrower, while the best-fit is little changed relative to \PrevRunPeriod.

\begin{figure*}[htb]
	\center
	\begin{subfigure}[b]{0.49\textwidth}
		\includegraphics[width=0.9\textwidth]{images/data/comparison/theta23_wrc_normal_1-9d_v_1-9c}
		\caption{\ssqthtwothree NH with reactor constraint}
	\end{subfigure}
	\begin{subfigure}[b]{0.49\textwidth}
		\includegraphics[width=0.9\textwidth]{images/data/comparison/theta23_wrc_inverted_1-9d_v_1-9c}
		\caption{\ssqthtwothree IH with reactor constraint}
	\end{subfigure}
	\begin{subfigure}[b]{0.49\textwidth}
		\includegraphics[width=0.9\textwidth]{images/data/comparison/theta23_worc_normal_1-9d_v_1-9c}
		\caption{\ssqthtwothree NH without reactor constraint}
	\end{subfigure}
	\begin{subfigure}[b]{0.49\textwidth}
		\includegraphics[width=0.9\textwidth]{images/data/comparison/theta23_worc_inverted_1-9d_v_1-9c}
		\caption{\ssqthtwothree IH without reactor constraint}
	\end{subfigure}
	\caption {
		The expected $\Delta\chi^2$ distribution as a function of \ssqthtwothree with and without reactor constraint. Normal and inverted hierarchy contours are independent. All 5 samples were used to produce these fits.
	}
	\label{fig:th23_data_compare}
\end{figure*}
\clearpage

\subsection{Results for \dmsq}

% TODO: Update this for run 10
The result for \dmsq with reactor constraint and T2K only is shown in fig.~\ref{fig:dmsq_data_compare}. %The \RunPeriod data favour higher values of \dmsq relative to \PrevRunPeriod, with a slightly narrower contour.

\begin{figure*}[htb]
	\center
	\begin{subfigure}[b]{0.49\textwidth}
		\includegraphics[width=0.9\textwidth]{images/data/comparison/dmsq_wrc_normal_1-9d_v_1-9c}
		\caption{\dmsq NH with reactor constraint}
	\end{subfigure}
	\begin{subfigure}[b]{0.49\textwidth}
		\includegraphics[width=0.9\textwidth]{images/data/comparison/dmsq_wrc_inverted_1-9d_v_1-9c}
		\caption{\dmsq IH with reactor constraint}
	\end{subfigure}
	\begin{subfigure}[b]{0.49\textwidth}
		\includegraphics[width=0.9\textwidth]{images/data/comparison/dmsq_worc_normal_1-9d_v_1-9c}
		\caption{\dmsq NH without reactor constraint}
	\end{subfigure}
	\begin{subfigure}[b]{0.49\textwidth}
		\includegraphics[width=0.9\textwidth]{images/data/comparison/dmsq_worc_inverted_1-9d_v_1-9c}
		\caption{\dmsq IH without reactor constraint}
	\end{subfigure}
	\caption {
		The expected $\Delta\chi^2$ distribution as a function of \dmsqtwothree with and without reactor constraint. Normal and inverted hierarchy contours are independent. All 5 samples were used to produce these fits.
	}
	\label{fig:dmsq_data_compare}
\end{figure*}

\clearpage